% =====================================================================
% Reviewer a3Zu -- Rating 5 (Accept), Confidence 3
% Owns Rebuttal Tables 7 and 8.
% Cross-references Rebuttal Tables 4, 9 (appendix.tex) and Tables 5, 6 (gbsa.tex).
% =====================================================================

\section*{Reviewer a3Zu}

\noindent\textit{Rating 5 (Accept), Confidence 3.}

\subsection*{Weaknesses}

\begin{itemize}
\item \textit{``In cognitive distortion analysis, authors scan reasoning traces for language
associated with loss chasing and similar patterns, but it is never validated against human
judgement. No human annotator checked a sample of the flagged outputs to confirm they actually
reflect the distortion they are supposed to capture. Despite this, the cognitive distortion
framing appears throughout the paper including in the abstract, carrying more rhetorical weight
than the underlying method can support.''}
\plan{키워드 측정치를 ``도구''에서 특이도 미확립 기준선으로 강등하고, 측정량을
gambling-related linguistic markers로 개명했다. 다만 초록에는 cognitive distortion이 등장하지
않으므로 철회 범위는 §2 구성개념 정의·§3 Finding 5·부록으로 정정했고, Finding 5가 이미 달고
있던 한정 문구를 인용해 원논문 주장이 이미 output-side로 제한돼 있었음을 밝혔다. 블라인드 인간
코딩 도구는 배포했으나 라벨은 아직 수집 전이다.}
\end{itemize}

\subsection*{Questions}

\begin{itemize}
\item \textit{``What is the parsing error or ambiguity rate for the moving-target metric? If a
meaningful share of outputs couldn't be reliably parsed, that uncertainty should be factored
into how much weight we put on the goal-escalation finding.''}
\plan{요청하신 모호성 비율은 측정하지 못했다고 먼저 밝히고, 비추출률은 분모와 함께 8월 3일까지
상한으로 보고한다. Figure 4(c)와 부록 표 12셀의 정확 재현을 확인한 뒤 계기 결함 2건(공개가중치
no-goal 셀은 구조적 0=측정불가; 무요청 목표 추출 64.4\%/53.3\%)을 공개하고, 엄격 정의(달성 후
상향)는 보정치가 아니라 민감도 분석으로만 제시한다: 계기가 대칭인 API 표본(조건당 1,600게임)에서
공표 지표 2.83$\times$ 대 엄격 2.24$\times$, 대비가 정의되는 4개 모델 전부 방향 일치(개방가중치
2종은 no-goal 대비 정의 불가). 남은 민감도 검사 한 건은 아직 측정 전이라 방향을 단정하지 않는다.}

\item \textit{``Have you looked at whether humans playing the exact same slot-machine game under
the same autonomy conditions show similar patterns?''}
\plan{인간 피험자 실험은 수행하지 않았음을 명시하고, 2023년 슬롯머신 발화 연구의 인간 코더
빈도를 외부 참조 분포로만 사용하되 정량적 대응을 함의하던 문장은 삭제했다.}

\item \textit{``All prompt conditions are zero-shot. What happens if you include one example of
someone playing cautiously and stopping at the right time? That single demonstration might be
enough to anchor the model and reduce the gambling-like behavior, which would suggest the effect
is at least partly about the model having nothing concrete to calibrate against, rather than
autonomy fundamentally breaking its decision-making. You could also try the opposite with one
example of escalating play under the BASE prompt with no autonomy framing to see if
demonstration alone is enough to produce the effect without any of the goal-setting or
reward-maximization modules.''}
\plan{DEMO-cautious와 DEMO-escalate 두 팔을 지시하신 그대로 사전등록했으나 아직 실행 전이며,
이미 확보된 합리성 지시 결과(Rebuttal Table 5, 6)는 방향이 일치하는 부분 증거로만 제시했다.}
\end{itemize}

\subsection*{Limitations and Formatting}

\begin{itemize}
\item \textit{``yes \textperiodcentered{} table 2 could be better formatted''}
\plan{Table 2를 열 그룹 정렬, 단위의 헤더 이동, 조건별 n을 캡션에서 표 본문으로 이동하는
방식으로 재조판했다.}
\end{itemize}

\redrule

\subsection*{\color{red!60!black}Author Response}

We thank you for a review in which every concern is stated as something we could act on, and for
the sentence that names our central problem most fairly: the framing carries ``more rhetorical
weight than the underlying method can support.'' We accept that concern, and the revision reports
language-frequency contrasts only until human validation is complete.

Your points ask one question of four different instruments --- whether the number we report can
carry the name we gave it: a keyword rule never checked against human judgement, a goal extractor
whose error rate was never stated, a paradigm with no human reference, and a zero-shot prompt with
nothing to calibrate against. Here is where each landed.

\begin{itemize}
\item \textbf{Was the cognitive-distortion measure validated against human judgement?} No, and
the gap is ours. We rename the measured quantity to \emph{gambling-related linguistic markers}
and withdraw the psychological claim from every place it is made --- §1's framing, §2's construct
definition, §3's opening and closing paragraphs, §3 Finding~5 and the appendix. One correction to
the premise: the submitted abstract does not use the term and does not report the measure, so it
needs no retraction. A blinded human-coding instrument is deployed, with its decision rule fixed
numerically before any label; labels follow by 3~August.
\item \textbf{What is the parsing error or ambiguity rate for the moving-target metric?} Not
measured, and we say so rather than substitute a narrower number for the one you asked for. We
reproduced the published Figure~4(c) values exactly, disclose two instrument defects your
question led us to find, and add a strict-definition sensitivity analysis under which the
goal-setting contrast is $2.24\times$ against the published metric's $2.83\times$ on the sample
where the instrument is symmetric. The non-extraction rate, with its denominator, follows by
3~August.
\item \textbf{Have humans played the same game under the same conditions?} No, and we cannot run
that study within a rebuttal. We use a 2023 human think-aloud slot-machine study as an external
reference distribution only, and the revision removes every sentence that implied a quantitative
correspondence.
\item \textbf{What happens with one cautious demonstration, and one escalating demonstration
under BASE?} Not yet run. Both arms are registered exactly as you specified them, seed-matched at
cap \$70 with $n = 100$ games per cell, and we report by 3~August; the completed
rationality-instruction factorial points the same way but cannot substitute for a demonstration.
\item \textbf{Table 2 formatting.} Fixed. Column groups align, units move into the header, and
the per-condition $n$ moves into the table body.
\end{itemize}

If any part of our response falls short, we would be glad to take it further during the
discussion period.

% ---------------------------------------------------------------------
\subsubsection*{\textbf{[Weakness 1]} The measure was never validated against human judgement}
% ---------------------------------------------------------------------

We understand your concern. As we understand it, you are pointing at a gap between what our
instrument measures and what its name claims: a keyword rule counts surface expressions, whereas
``cognitive distortion'' is a psychological attribution that only a human judgement can license.

\textbf{Verdict: the gap is accepted; one part of the premise needs correcting.} No human
annotator checked the flagged outputs before submission. That gap is real and it is ours.
\textbf{Thank you for this sharp observation. It led us to build and deploy a blinded human-coding
instrument, and to re-run the language analysis under five alternative instruments.}

\textbf{The correction: Finding 5 already scoped the claim.} The label does not appear in the
abstract, as noted above.
Finding 5, which does report the measure, already states that we ``treat these as output-side interpretive
frames rather than causal mechanisms,'' that the analysis ``is not evidence that the model
independently discovers those distortions, only that high-risk regimes are accompanied by
loss-recovery and control-like justifications in the generated reasoning,'' and that these traces
``do not validate a clinical diagnosis or a causal reasoning mechanism.''
However, we think your point survives that hedge.
A qualification inside one paragraph does not offset a construct name used as the name of the
measure itself.
We therefore change the name, not only the hedge.

\textbf{Three kinds of validity are at stake.} Each kind receives a different answer.
We separate them so that none borrows support from another.

\begin{itemize}
\item \textbf{Frequency validity} asks whether language occurs more often in one condition than another.
The method measures this quantity, and the revised paper claims only this contrast.
The method passes this test. Rebuttal Table 7 reports the result.
\item \textbf{Content validity} asks whether flagged text expresses the construct named by its label.
This is the validity you requested. Only human coding can establish it.
We built and blinded the instrument, but collected no labels.
\item \textbf{Criterion validity} asks whether the measure agrees with an independent instrument targeting the same construct.
We attempted that test with a five-instrument battery, but the battery cannot serve as one.
The instruments are not independent, so their agreement is not criterion evidence.
Rebuttal Table 7 must be read with that limit attached.
\end{itemize}

\textbf{The frequency contrast persists across instruments.} This re-analysis asks whether the
contrast depends on the particular word list we happened to write. We varied the word list and held
the corpus fixed, rerunning all 19,200 games and 190,300 decisions under five instruments. The
contrast stays positive in 6 of 6 models under every one of them.

An instrument is one list of regular expressions grouped into named frames; a frame is the
expression group for one construct, such as loss-chasing. The five instruments are:

\begin{enumerate}[leftmargin=1.8em]
\item the paper's own four frames;
\item a frozen codebook drawn from constructs that converge across validated gambling instruments;
\item that codebook with every goal-coupled expression ablated;
\item expressions modelled on the Gambling Related Cognitions Scale (GRCS);
\item think-aloud-style expressions.
\end{enumerate}

\noindent
We ran four of these --- the paper's frames, the frozen codebook, the GRCS-style set and the
think-aloud set --- twice each, once as written and once with a polarity correction that drops
sentences where the model rejects rather than expresses the cognition. That gives the eight
variants quoted throughout. The goal-coupled-ablated codebook is reported separately, as the
ablated column of Rebuttal Table 7. Throughout, pp means percentage points.

Rebuttal Table 7 gives the goal minus no-goal contrast per model under three of these instruments.

\vspace{0.5em}
\noindent\textbf{Rebuttal Table 7.} Goal minus no-goal contrast in pp, per model, under three
instruments, full corpus of 19,200 games and 190,300 decisions. Positive in 6 of 6 models under
every instrument, and under every one of the eight variants.

\begin{center}
\footnotesize
\begin{tabular}{lccc}
\toprule
\textbf{Model} & \textbf{Original, window-scoped} & \textbf{Convergent codebook} & \textbf{Goal-coupled ablated} \\
\midrule
GPT-4o-mini      & $+77.9$ & $+37.2$ & $+29.6$ \\
GPT-4.1-mini     & $+75.1$ & $+44.9$ & $+41.9$ \\
Gemini-2.5-Flash & $+77.9$ & $+42.2$ & $+41.7$ \\
Claude-3.5-Haiku & $+30.4$ & $+18.6$ & $+17.2$ \\
LLaMA-3.1-8B     & $+18.1$ & $\phantom{0}+4.2$ & $\phantom{0}+3.1$ \\
Gemma-2-9B       & $+38.8$ & $+16.0$ & $+12.7$ \\
\bottomrule
\end{tabular}
\end{center}

\vspace{0.4em}
\noindent
The autonomy contrast is weaker and we report it as such.
Variable minus fixed is positive in 5 of 6 models in seven of the eight variants, and in 4 of 6 in
the raw original variant, where GPT-4o-mini is $-0.1$ pp.
Gemini is negative in every variant, at $-6.1$ to $-11.9$ pp.
We attribute Gemini's sign to the model rather than to an instrument defect.

\textbf{We restored a scoping step we had earlier dropped, and it moved the numbers in our
favour.} The paper's analysis restricts the pattern-belief frame to conditions whose prompt does
not mention hidden patterns, and the loss-chasing frame to decisions that follow a loss.
An earlier version of this battery omitted that restriction and reported a smaller goal contrast.
The hidden-patterns module is crossed with the goal module, so its prompt text raised the
pattern-belief hit rate in the goal and no-goal arms alike: it diluted the contrast rather than
inflating it.
The scoped column in Rebuttal Table 7 is the one that reproduces the paper's statistic, and it is
larger in every model.
The convergent and ablated columns are unaffected, because neither contains a pattern-belief or
loss-chasing frame.

\textbf{The ablated column is the strongest check.} In the original instrument, the
\texttt{goal\_escalation} frame dominates the goal contrast at $+65$ to $+95$ pp.
That is close to circular, because the goal module creates the goals that can escalate.
The frozen codebook has no goal frame; its contrast comes from \texttt{illusion\_of\_control}
($+16.7$ to $+58.4$ pp) and \texttt{impaired\_control} ($+13.5$ to $+50.6$ pp), and it survives
the removal of every goal-coupled expression at 6 of 6 models.
LLaMA is the weakest case: 87.6\% against 84.5\%, a $+3.1$ pp gap with $z = 2.50$ and
$p \approx 0.012$, with overlapping Wilson intervals.
We report LLaMA as a pass, but not a comfortable one.

\textbf{Three limits apply, and we raise all three ourselves.} The first two prevent the battery
from serving as a criterion test.

\begin{itemize}
\item \textbf{The instruments are not independent.} Several expressions appear verbatim in more than one instrument.
Shared wording prevents us from treating agreement between lists as convergence.
\item \textbf{Two instruments were written after the original result was known.} We wrote the GRCS-style and think-aloud-style expressions during the rebuttal period.
They are robustness checks, not a pre-registered replication.
\item \textbf{One frame partly measures prompt echo.} Prompt echo means repeating wording supplied by the prompt.
The hidden-patterns module uses the prompt text ``This slot machine may have hidden patterns,''
and the pattern and hidden expressions match that text directly.
This is why the scoping step above excludes those conditions from the pattern frame.
\end{itemize}

\textbf{We can identify two false positives without human labels.} A \texttt{self\_serving\_bias}
expression labels ``stopping now is the smart decision'' as a distortion, when that sentence
describes a rational refusal.
\texttt{illusion\_of\_control} also misfires in the variable arm, where stake size genuinely remains under the model's control.
Appendix~A.3 prints both expressions in full, and Appendix~A prints every expression together with
the frozen file's SHA-256.

\textbf{We searched for a validated public lexicon for this task and found none.} Rebuttal Table 8 traces the sources.
It shows what our four constructs inherit and what we wrote ourselves.

\vspace{0.5em}
\noindent\textbf{Rebuttal Table 8.} Prior-work lineage of the four constructs: what each source
supplies, in what form, and whether it can serve as a free-text lexicon.

\begin{center}
\footnotesize
\begin{tabular}{p{0.185\linewidth}p{0.305\linewidth}p{0.145\linewidth}p{0.215\linewidth}}
\toprule
\textbf{Source} & \textbf{What it supplies} & \textbf{Form} & \textbf{Free-text lexicon?} \\
\midrule
Goodie \& Fortune (2013), \textit{Psychology of Addictive Behaviors} 27(3), 730--743 &
Convergence across validated instruments: all prominent instruments include illusion of control,
almost all include gambler's fallacy &
Review of instruments &
No. It is our source for which constructs to use, not for any wording \\
\addlinespace
Raylu \& Oei (2004), GRCS &
Five subscales: inability to stop, interpretative bias, illusion of control, gambling
expectancies, predictive control &
Self-report questionnaire items &
No. Items are answered by a person, not matched against text \\
\addlinespace
Toneatto (1999) typology &
Magnification of skill, minimisation of others' skill, superstitions, interpretive biases,
temporal telescoping, selective memory, predictive skill, illusion of control over luck,
illusory correlation, entitlement, omnipotence, magical thinking &
Clinical typology &
No. It names constructs and supplies no matching rule \\
\addlinespace
2023 simulated slot-machine verbalisation study &
Eight coded categories, with frequencies gambler's fallacy 57, near-miss 47, illusion of
control 46 &
Human think-aloud coding scheme &
No, but it is a human reference distribution in a related paradigm \\
\addlinespace
Bathina et al. (2021), \textit{Nature Human Behaviour} &
12 categories, 241 n-grams &
Lexicon &
It is a lexicon, but for general cognitive distortion, not gambling-specific ones \\
\addlinespace
Smith et al., \textit{PLOS Digital Health} &
DSM-5 plus GRCS annotation guide for problem-gambling content &
Manual annotation guide &
No. It is the closest gambling-specific public resource and it is a guide, not a lexicon \\
\bottomrule
\end{tabular}
\end{center}

\vspace{0.4em}
\noindent
No row supplies a lexicon we could use. We therefore wrote the expressions ourselves from the four
definitions, then froze the file and recorded its hash before computing any statistic.

\textbf{We define a new target quantity and state when it is honest.} The method measures one quantity.
It measures the rate at which gambling-related language appears in a reasoning trace.
We rename this quantity from \textbf{cognitive distortions} to \textbf{gambling-related
linguistic markers}.

The rename changes the measurement label, not the codebook contents.
The surviving contrast still concerns constructs that Goodie and Fortune find convergent.

The rename is honest only if the psychological claim goes from every place it was made.
It was made in §1's framing, in §2's construct definition, in §3's opening and closing paragraphs,
in §3 Finding 5 and in the appendix, and it is withdrawn in all of them.
The abstract needs no retraction, because it never carried the claim.
We do not claim that this language reflects the model's cognition, and we do not claim that it
predicts the next decision.

\textbf{We built and blinded the human-coding instrument. It remains a pilot, not a completed
validation. This experiment is under way.} The instrument contains 100 items: 25 for each of the
four frozen constructs, and within each construct 12 regex-flagged and 13 unflagged items, so that
missed genuine instances are estimated rather than assumed.
Responses of exactly 500 characters are excluded, that length being the signature of a storage
truncation we disclose in Appendix~B3.
Coders see only the trace and the construct name; model, condition, flag status and matched span
are withheld.
Items are drawn from the six models in strict rotation, the presentation order is shuffled, and
the seed is fixed at 24231 so the draw is reproducible.

Three coders will label the items: two authors and one person outside this project.
We report author and non-author contrasts separately, and compute $\kappa$, the inter-coder
agreement statistic, before adjudication.
We fixed the decision rule numerically before seeing any label.
The rule has three clauses.

\begin{itemize}
\item $\kappa$ below 0.60: we make no quantitative statement at all.
\item A 95\% lower bound on a frame's precision below 0.50: we remove that frame's numbers from the paper's body.
\item A 95\% interval on the human-labelled variable-minus-fixed contrast that includes zero: we withdraw that claim, which the paper reports as Finding 5.
\end{itemize}

Recall has no decision gate; we report it descriptively.
\textbf{We have collected no labels, and we have not yet recruited the non-author coder.}
We are not asking you to credit a result. We are showing you the instrument and the rule that will
judge it. We will complete it and report the outcome by 3 August, regardless of whether it favours
us. We intend to reflect this in the camera-ready.

% ---------------------------------------------------------------------
\subsubsection*{\textbf{[Question 1]} Parsing error or ambiguity rate for the moving-target metric}
% ---------------------------------------------------------------------

We understand your concern. As we understand it, you are pointing at a metric whose input is a
free-text extraction: if the extractor misfires, the moving-target rate inherits that error, and
the paper never reported how often it misfires.

\textbf{We cannot give you the rate you asked for.} We have not measured the extractor's false-positive rate.
We have not measured how often an extracted goal value is semantically ambiguous.
Your question asks for more than our current measurements provide.
Reporting a narrower number as your requested rate would misrepresent the evidence.
What we can do is verify the pipeline, disclose the two instrument defects your question led us
to find, and bound the finding with a sensitivity analysis. We do the three in that order.

\textbf{The published numbers reproduce exactly.} Re-running the figure's own extractor on the
figure's own corpus (9,600 games, 2,400 per condition) returns the published Figure~4(c) values
to one decimal --- BASE 17.0, M 11.0, G 49.8, GM 47.8 --- and all twelve cells of the appendix
table's moving-target column reproduce cell for cell. Everything below therefore concerns the
paper's own pipeline, not a reconstruction of it.

\textbf{Your question also surfaced a description error we have since corrected, and we disclose
it here rather than let you find it.} The submitted §2 defined the moving-target rate as the
fraction of games in which the model \emph{raises its self-set goal after meeting it}. The
implementation never tested whether the goal had been met; it flags any upward revision during
play. The definition and the code therefore measured different quantities, and the code is the one
that produced every number in the paper.

We changed the definition to match the implementation --- ``revises its self-set goal upward
mid-game, with or without having reached it'' --- rather than change the code, because rewriting
the code would invalidate the reported figures without new data. That fix was incomplete, and we
disclose the remainder: the clinical-framing paragraph earlier in §2 still describes the
construct as achievement-conditional in two phrases, ``shifting one's own goals upward after
meeting them'' and ``raising the target once it has been reached.'' Both receive the same
correction in the camera-ready, so that no sentence in the paper claims an achievement test the
metric does not perform.

\textbf{Defect 1: the open-weight no-goal cells are structural zeros.} For LLaMA and Gemma the
pipeline reads stored goal fields rather than extracting from text, and those fields are written
only when the goal prompt is on: coverage is 368--396 of 400 games per cell under G and GM, and
0 of 400 under BASE and M for both models. The two arms therefore fail differently. The goal
arm has partial coverage, and a missing field can only hide an escalation, so that side is
conservative; the no-goal arm has no coverage at all, so its contrast is not defined --- a zero
there means \emph{not measurable}, not ``no escalation occurred.'' The pooled BASE and M figures are therefore the API rate diluted by 800
structural zeros per condition, and the dilution is exact: the API-only rates 25.4 and 16.5
become 16.9 and 11.0 after multiplying by $2/3$, against the published 17.0 and 11.0. Because
the goal arms are not diluted, the asymmetry inflates the published contrast from $2.83\times$
on the comparable API sample to $3.5\times$ after pooling. The camera-ready marks those two
cells \emph{not measurable} instead of 0.0\%, and restates the appendix sentence coupling
bankruptcy with the moving-target rate so that it claims the four models in which the contrast
is defined, not six.

\textbf{Defect 2: goals are extracted where none were requested.} In the API no-goal cells the
extractor returns a goal in 64.4\% of BASE games and 53.3\% of M games, although no goal was
asked for. We do not know what those matches are. One reading is that the extractor is matching
a balance the game has already passed, which would fit the near-coincidence of the published and
strict metrics in those cells (25.4 vs 24.6; 16.5 vs 14.8); that reading is a hypothesis, not a
measured fact, and bet amounts, loss limits, cash-out thresholds, goals a model volunteers
unprompted, or ordinary numeric prose could produce the same matches. A structural point holds
either way: the denominator is all games in the condition, so the metric mixes a change in
behaviour with a change in how often a goal is stated at all --- 97.5\% of G games state one,
against 43.0\% of BASE games.

\textbf{A sensitivity analysis under the submitted definition. It is a sensitivity analysis, not
a corrected value.} Recomputing with the achievement test added and nothing else changed:

\begin{center}
\footnotesize
\begin{tabular}{lccc}
\toprule
\textbf{Condition} & $n$ & \textbf{Published metric} & \textbf{Strict (raise after reaching)} \\
\midrule
BASE & 2,400 & 17.0 [15.5, 18.5] & 16.4 [14.9, 17.9] \\
M    & 2,400 & 11.0 [9.8, 12.3]  & \phantom{0}9.9 [8.7, 11.1] \\
G    & 2,400 & 49.8 [47.8, 51.7] & 34.4 [32.5, 36.3] \\
GM   & 2,400 & 47.8 [45.8, 49.8] & 34.0 [32.2, 36.0] \\
\bottomrule
\end{tabular}
\end{center}

\noindent
Because the pooled control arm contains the unmeasurable open-weight cells, the ratio should be
read on the API-only sample, where the instrument is symmetric: strict BASE 24.6, M 14.8, G
46.1, GM 42.2 ($n = 1{,}600$ per condition), a contrast of $2.24\times$ against $2.83\times$
for the published metric on the same sample. Per model under the strict rule (goal versus
no-goal arm, $n = 800$ each): GPT-4o-mini 61.4 vs 35.2, GPT-4.1-mini 45.8 vs 2.1, Gemini 38.6
vs 15.4, Claude 30.8 vs 26.0 --- positive in 4 of 4 models in which the comparison is defined,
Claude the narrowest; the two open-weight models have no defined comparison. The defensible
claim is that goal-setting raises the observed rate of upward goal revision by roughly
$2.2\times$ to $2.8\times$ on the sample where the instrument is symmetric, and that the
direction holds in every model in which the contrast can be computed. We no longer write
``roughly triples'' without naming the pooled instrument that produces it.

\textbf{Five objections this analysis invites, all of which we concede in advance.} The strict
rule uses the same extractor and inherits its errors. The stricter definition arrives after the
definition was already relaxed once, which is why we present it only as a sensitivity analysis.
The pooled control arm is partly unmeasured. The denominator conflates behaviour with goal
observability. And the goal prompt itself elicits more numeric talk, so higher extraction in the
goal arm is partly mechanical.

\textbf{We can report the following scope precisely.} The moving-target metric flags a game when
its extracted goal values rise during play. Our protocol does not require the model to restate
the goal each round, so a decision that yields no goal value is not necessarily a parsing
failure. We will report the non-extraction rate, with its denominator, by 3~August, and the
revision states it as an upper bound on parsing trouble at the point of reporting.

\textbf{Missingness biases the rate downward; mis-extraction need not.} A missing goal value can
only delete entries from the observed sequence, and deleting entries from a sequence that never
rises cannot make it rise, so non-extraction hides upward revisions rather than manufacturing
them. Mis-extraction --- Defect 2 --- can bias in either direction, which is why the finding
rests on the strict rule and the symmetric API sample rather than on this argument alone.

\textbf{One registered check is still running, and we will not pre-announce its direction.} It
restricts the metric to games with at least two observed goal values, which removes the games
where missingness is worst. We will report it with its denominator by 3~August.

\textbf{We correct one description of the extractor.} The body implied a narrow
\texttt{goal/target \$N} regular expression. The implementation is broader: nine patterns applied
in order to the lower-cased response, of which the first to yield a value between \$50 and
\$10,000 wins, taking the last match within that pattern. Three of the nine require no goal word
at all --- \texttt{aim for/to \$N}, a single pattern covering both \texttt{reach \$N} and
\texttt{get to \$N}, and \texttt{balance ... of at least \$N}. The remaining six key on
\texttt{goal} or \texttt{target}, including bare \texttt{\$N goal} and a permissive
\texttt{goal: ... N dollars}.
The narrow expression measures a different quantity, and the revision describes the implementation
accurately.

\textbf{We found a second parsing defect during the rebuttal.} The bet/stop parser, which
classifies betting versus stopping, takes the first ``final decision'' match in a response and
tests betting tokens before stopping tokens. It therefore records a wager when a response
discusses betting before ending with a stop.
We re-parsed every stored response under a corrected rule that uses the last match and classifies
on the first word inside it. Rebuttal Table 4 in the appendix gives the flip counts, the
adjudicable denominators and the resulting rates for all three corpora.

The flips concentrate in one model. Every flip in the re-run corpus belongs to Claude-Haiku, so
that model's fixed-arm participation counts are essentially all misparsed stops, and they never
appear here unannotated.

\textbf{The model identity also requires correction.} The frozen configuration declares one
model, and every run used another:

\begin{center}
\small declared \texttt{claude-3-5-haiku-20241022}; \; ran \texttt{claude-haiku-4-5-20251001}
\end{center}

\noindent
The revision reports the model we actually ran in Appendix~B7.

\textbf{The moving-target metric uses a different extractor, so this defect does not affect it.}

% ---------------------------------------------------------------------
\subsubsection*{\textbf{[Question 2]} Humans playing the same game under the same conditions}
% ---------------------------------------------------------------------

We understand your concern. As we understand it, you are asking whether the pattern we attribute to
autonomy is specific to models at all, since a human given the same freedoms in the same game might
produce the same escalation.

\textbf{Verdict: no, we did not run it, and we cannot run it within this rebuttal.} We ran no human
participants on this task and did not seek the ethics review such a study requires. A
within-rebuttal human-subjects study on a gambling paradigm would be inappropriate regardless of
timing.

\textbf{We can offer an external anchor instead of a blanket disclaimer.} The 2023 simulated
slot-machine verbalisation study reports human coder frequencies for eight categories, including
gambler's fallacy 57, near-miss 47 and illusion of control 46; Rebuttal Table 8 places it in the
source record. We add those eight categories as a second coding layer alongside the four frozen
constructs, leaving the frozen instrument unchanged so the comparison stays auditable.

\textbf{The comparison has a strict bound.} That study collected human verbalisations under a
different protocol, bankroll and cap structure. It is a reference distribution for a related
paradigm, not a matched control for ours. We claim no quantitative correspondence between our
models and human players, and the revision removes every sentence that implied one. We intend to
reflect this in the camera-ready.

% ---------------------------------------------------------------------
\subsubsection*{\textbf{[Question 3]} One cautious demonstration, and one escalating demonstration}
% ---------------------------------------------------------------------

We understand your concern. As we understand it, you are proposing a rival account of our result:
the model may escalate not because autonomy breaks its decision-making, but because a zero-shot
prompt gives it nothing concrete to calibrate against --- and you propose the two demonstrations
that would separate the accounts.

\textbf{Thank you for this sharp observation. It led us to register both arms exactly as you
specified.} \textbf{Verdict: accepted in both directions, not yet run.} Your premise is correct:
every prompt condition in the submitted paper is zero-shot, with no worked example anywhere.

DEMO-cautious prepends one worked example of a player who stops early, under the autonomy
conditions; on your account risk-taking should fall sharply there. DEMO-escalate prepends one
worked example of escalating play under the base prompt with no autonomy framing; on your account
the effect should appear without an autonomy module. Both run at cap \$70 with $n = 100$ games per
cell, reusing the framing factorial's random seeds so the games are matched across arms.
We pre-register the direction of both predictions and deliberately register no numeric gate: one
demonstration and two competing accounts give no principled threshold to fix in advance, and
declaring none now prevents adjusting one later. Rebuttal Table 9 in the appendix records both
arms and the decision rule.

\textbf{This experiment is under way. We will complete it and report by 3 August. We intend to
reflect this in the camera-ready.}

\textbf{One completed result bears on your hypothesis without settling it.} The rationality
instruction in the completed framing factorial nearly eliminates participation in the models that
gamble most. Reviewer gbSA's Question 3 contains our full answer, including the cells where the
instruction is not a complete off switch; we summarise the conclusion here. Rebuttal Table 6 in
that section gives the main effect on variable-arm participation per model, and Rebuttal Table 5
the underlying counts for the 32 completed cells. Gemini is the cleanest of the three, because its
runner passes no system message, whereas GPT-4.1-mini and GPT-4o-mini carry standing system
messages calling the decision maker ``cautious, rational'' and ``rational'' respectively, and our
factorial inserts factor text into the user prompt only (Appendix~B12).

However, we think that result cannot stand in for the arms you asked for. The factor is an
instruction, not a demonstration, so the design cannot separate information about the odds from
the endorsement of stopping. The demonstration arms make exactly that separation, which is why we
adopted them as specified.

% ---------------------------------------------------------------------
\subsubsection*{\textbf{[Formatting]} Table 2}
% ---------------------------------------------------------------------

We understand your concern. As we understand it, Table 2 is hard to read at a glance because the
column groups do not line up and the reader must go to the caption for the units and the
per-condition $n$.

\textbf{Thank you for this sharp observation. It led us to re-typeset the table.} Column groups
now align, units move into the header, and the per-condition $n$ moves from the caption into the
table body. We intend to reflect this in the camera-ready.
